% GENERATED FILE. Edit canonical lessons, then run npm run generate:books.
% canonical-source-hash: fnv1a64:5af0fa86996f6de7
% canonical-lessons: FR-C487-amoureux, FR-C487-coup-de-foudre, FR-C487-draguer, FR-C487-tendresse, FR-C487-caresser

\chapter{Amoureux}
\label{ch:fr-amoureux}
\hlchaptermodality{487}

\begin{chapteropening}
\textbf{By the end of this chapter:} \emph{I can talk about falling in love.}

Complete the last lesson of chapter 487: I can talk about falling in love.
\end{chapteropening}

\section[amoureux]{amoureux — in love; lover}

\label{lesson:FR-C487-amoureux}



\begin{quote}
Before the new one: say the French for to yawn, then the French for nap.
\end{quote}

\subsection*{You'll want to know: amoureux}
\textbf{amoureux} — \textquotedblleft{}in love; lover\textquotedblright{}.

\begin{quote}
Ils sont amoureux depuis le lycée.
\end{quote}

\begin{quote}
\emph{They have been in love since secondary school.}
\end{quote}

Feminine \emph{amoureuse}. \emph{Tomber amoureux de quelqu'un} — to fall in love with someone.

\begin{cousinweb}
\emph{amoureux} is inherited from Old French \emph{amoros} (EtymDB), built on \emph{amour}, from Latin \emph{amor} \textquotedblleft{}love\textquotedblright{} (Lewis \& Short).
\end{cousinweb}

\begin{grammarlens}[title={What you've built}]
One more word for telling, explaining and describing at B1.
\end{grammarlens}

\subsection*{Your turn}
\begin{itemize}
\raggedright
  \item \emph{Say it:} \emph{amoureux}
  \item \emph{Say it:} \emph{amoureux}, once more
  \item \emph{Say it:} a sentence of your own with \emph{amoureux}
  \item \emph{Recall:} say the French for to yawn, then the French for nap, then say \emph{amoureux} again
\end{itemize}

\begin{tcolorbox}[breakable,colback=teal!4,colframe=teal!35!black,title={Before you move on}]
What does amoureux mean? (\textquotedblleft{}In love; lover\textquotedblright{}.) Say it once more.
\end{tcolorbox}
\section[le coup de foudre]{le coup de foudre — love at first sight}

\label{lesson:FR-C487-coup-de-foudre}



\begin{quote}
Before the new one: say the French for nap, then the French for in love.
\end{quote}

\subsection*{You'll want to know: le coup de foudre}
\textbf{le coup de foudre} — \textquotedblleft{}love at first sight\textquotedblright{}.

\begin{quote}
Quand ils se sont rencontrés, ça a été le coup de foudre.
\end{quote}

\begin{quote}
\emph{When they met, it was love at first sight.}
\end{quote}

Literally a lightning strike. \emph{Avoir un coup de foudre pour} also works for places or objects.

\begin{cousinweb}
\emph{le coup de foudre} is \textquotedblleft{}a stroke of lightning\textquotedblright{}: \emph{coup} comes from Latin \emph{colaphus} \textquotedblleft{}a punch\textquotedblright{}, and \emph{foudre} is inherited from Latin \emph{fulgur} \textquotedblleft{}lightning\textquotedblright{} (Lewis \& Short).
\end{cousinweb}

\begin{grammarlens}[title={What you've built}]
One more word for telling, explaining and describing at B1.
\end{grammarlens}

\subsection*{Your turn}
\begin{itemize}
\raggedright
  \item \emph{Say it:} \emph{le coup de foudre}
  \item \emph{Say it:} \emph{le coup de foudre}, once more
  \item \emph{Say it:} a sentence of your own with \emph{le coup de foudre}
  \item \emph{Recall:} say the French for nap, then the French for in love, then say \emph{le coup de foudre} again
\end{itemize}

\begin{tcolorbox}[breakable,colback=teal!4,colframe=teal!35!black,title={Before you move on}]
What does le coup de foudre mean? (\textquotedblleft{}Love at first sight\textquotedblright{}.) Say it once more.
\end{tcolorbox}
\section[draguer]{draguer — to chat up, to flirt with (informal); to dredge}

\label{lesson:FR-C487-draguer}



\begin{quote}
Before the new one: say the French for in love, then the French for love at first sight.
\end{quote}

\subsection*{You'll want to know: draguer}
\textbf{draguer} — \textquotedblleft{}to chat up, to flirt with (informal); to dredge\textquotedblright{}.

\begin{quote}
Au café, un garçon a essayé de la draguer.
\end{quote}

\begin{quote}
\emph{At the café a boy tried to chat her up.}
\end{quote}

Informal. In Québec people say \emph{cruiser}. The first sense is to dredge a river.

\begin{cousinweb}
\emph{draguer} is built on \emph{drague} \textquotedblleft{}a dredge\textquotedblright{}, which was borrowed from English \emph{drag}. The sense \textquotedblleft{}to chat up\textquotedblright{}, like sweeping the streets, is recent and informal.
\end{cousinweb}

\begin{grammarlens}[title={What you've built}]
One more word for telling, explaining and describing at B1.
\end{grammarlens}

\subsection*{Your turn}
\begin{itemize}
\raggedright
  \item \emph{Say it:} \emph{draguer}
  \item \emph{Say it:} \emph{draguer}, once more
  \item \emph{Say it:} a sentence of your own with \emph{draguer}
  \item \emph{Recall:} say the French for in love, then the French for love at first sight, then say \emph{draguer} again
\end{itemize}

\begin{tcolorbox}[breakable,colback=teal!4,colframe=teal!35!black,title={Before you move on}]
What does draguer mean? (\textquotedblleft{}To chat up, to flirt with (informal); to dredge\textquotedblright{}.) Say it once more.
\end{tcolorbox}
\section[la tendresse]{la tendresse — tenderness, affection}

\label{lesson:FR-C487-tendresse}



\begin{quote}
Before the new one: say the French for love at first sight, then the French for to chat up.
\end{quote}

\subsection*{You'll want to know: la tendresse}
\textbf{la tendresse} — \textquotedblleft{}tenderness, affection\textquotedblright{}.

\begin{quote}
Elle regarde son fils avec tendresse.
\end{quote}

\begin{quote}
\emph{She looks at her son with tenderness.}
\end{quote}

Built on \emph{tendre} \textquotedblleft{}tender, soft\textquotedblright{}.

\begin{cousinweb}
\emph{la tendresse} is built on \emph{tendre} \textquotedblleft{}tender\textquotedblright{} with \emph{-esse} (EtymDB); \emph{tendre} is inherited from Latin \emph{tener} \textquotedblleft{}soft, delicate\textquotedblright{} (Lewis \& Short).
\end{cousinweb}

\begin{grammarlens}[title={What you've built}]
One more word for telling, explaining and describing at B1.
\end{grammarlens}

\subsection*{Your turn}
\begin{itemize}
\raggedright
  \item \emph{Say it:} \emph{la tendresse}
  \item \emph{Say it:} \emph{la tendresse}, once more
  \item \emph{Say it:} a sentence of your own with \emph{la tendresse}
  \item \emph{Recall:} say the French for love at first sight, then the French for to chat up, then say \emph{la tendresse} again
\end{itemize}

\begin{tcolorbox}[breakable,colback=teal!4,colframe=teal!35!black,title={Before you move on}]
What does la tendresse mean? (\textquotedblleft{}Tenderness, affection\textquotedblright{}.) Say it once more.
\end{tcolorbox}
\section[caresser]{caresser — to stroke, to caress}

\label{lesson:FR-C487-caresser}



\begin{quote}
Before the new one: say the French for to chat up, then the French for tenderness.
\end{quote}

\subsection*{You'll want to know: caresser}
\textbf{caresser} — \textquotedblleft{}to stroke, to caress\textquotedblright{}.

\begin{quote}
Le garçon caresse doucement le chat.
\end{quote}

\begin{quote}
\emph{The boy gently strokes the cat.}
\end{quote}

\emph{Une caresse} — a caress, a stroke.

\begin{cousinweb}
\emph{caresser} was borrowed from Italian \emph{carezzare} (EtymDB), built on \emph{carezza} \textquotedblleft{}a caress\textquotedblright{}, from \emph{caro} \textquotedblleft{}dear\textquotedblright{}, inherited from Latin \emph{cārus} \textquotedblleft{}dear\textquotedblright{} (Lewis \& Short).
\end{cousinweb}

\begin{grammarlens}[title={What you've built}]
One more word for telling, explaining and describing at B1.
\end{grammarlens}

\subsection*{Your turn}
\begin{itemize}
\raggedright
  \item \emph{Say it:} \emph{caresser}
  \item \emph{Say it:} \emph{caresser}, once more
  \item \emph{Say it:} a sentence of your own with \emph{caresser}
  \item \emph{Recall:} say the French for to chat up, then the French for tenderness, then say \emph{caresser} again
\end{itemize}

\begin{tcolorbox}[breakable,colback=teal!4,colframe=teal!35!black,title={Before you move on}]
What does caresser mean? (\textquotedblleft{}To stroke, to caress\textquotedblright{}.) Say it once more.
\end{tcolorbox}
